\chapter{Detektoreffizienz}

\todo{irgendwo erklären warum chi2 bei gradenfit doof ist}

Die Detektoreffizienz ist definiert als 
\begin{align}
	\epsilon = \frac{N_{E_\gamma \, \text{Photopeak}}}{N_{E_\gamma \, Detektor}}
\end{align}

wobei $N_{E_\gamma \, \text{Photopeak}}$ die Anzahl der im Photopeak einer Linie nachgewiesenen Photonen ist und $N_{E_\gamma \, \text{Detektor}}$ die Anzahl der Photonen dieser Linie ist, welche den Detektor erreichen\cite[6]{521versuchsanleitung}. 

Die Berechnung von $N_{E_\gamma \, \text{Detektor}}$ folgt aus der Aktivität der Probe sowie geometrischen Überlegungen; es wird angenommen, dass die Probe isotrop in alle Raumrichtungen strahlt. Dann ist der Anteil der Strahlung, welche den Detektor erreicht genau der Anteil einer Kugelfläche, welche der Detektor abdeckt. Bei runden Detektoren, zu denen die Probe normal und mittig platziert ist (was zumindest näherungsweise gegeben ist), ergibt sich die Fläche (als Teil einer Kugelschale), welche der Detektor verdeckt als

\begin{align}
	A_\text{Detektor} &= \int_0^\alpha \int_0^{2 \pi} \d{\theta} \d{\phi} r^2 \sin(\theta) = 2 \pi r^2 (1 - \cos(\alpha))\\
	\imply q &= \frac{A_\text{Detektor}}{A_\text{Kugel}} = \frac12 (1 - \cos(\alpha))
\end{align}

Der Winkel $\alpha$ ist hierbei der Öffnungswinkel eines Kegels, welcher von der Probe und dem Rand des Detektors aufgespannt wird, also $\alpha = \arctan(r_\text{Detektor} / d_\text{Probe})$ ($d_\text{Probe}$ ist dabei der Abstand der Probe zur Dektorvorderseite).
\todo{skizze davon in den anhang maybe}
Damit ist dann die Fläche gegeben als

\begin{align}
	q = \frac12 (1 - \frac{1}{\sqrt{1 + (\frac{r_\text{Detektor}}{d_\text{Probe}})}})
\end{align}


Die Aktivität der Proben zu einem Referenzzeitpunkt ist bekannt und muss nur mittels der Halbwertszeiten für den Versuchstag berechnet werden.
Dafür wird das Zerfallsgesetz genutzt \todo{source}. Während der relativ kurzen Messung wird die Aktivität als konstant betrachtet.
\begin{align}
	A = A_0 e^{- t \frac{\ln(2)}{T_{1/2}}}
\end{align}

Die Aktivitäten, gemessen am 01.10.2021, sind auf den Proben angegeben und zusammen mit den aktivitäten am Versuchstag (29.04.2026) in Tabelle \ref{tab:activity} zu finden. Hierbei wurde als Unsicherheit bei der Zeit ein Tag angenommen, da keine Uhrzeit der Messung bekannt ist. weiterhin wird auf die angegebene Aktivität ein Fehler von 5\% des angegebenen Wertes angenommen, da keine weitere Angabe vorliegt. Diese Unsicherheit mag natürlich völlig falsch sein, vermittelt bei den resultierenden Rechnungen aber ein Gefühl für die zu erwartende Genauigkeit.


\begin{table}[h]
    \centering
    \begin{tabular}{|c|c|c|c|} \hline
	Isotop & $T_{1/2}$ / a & A_0 / Bq & A(t_\text{versuch}) / Bq \\
	\hline
	137 Cs & \num{30.08(9)} & \num{4.05(20)e+05} & \num{3.64(18)e+05} \\
	152 Eu & \num{13.517} & \num{7.09(35)e+05} & \num{5.61(28)e+05} \\
	60 Co & \num{5.275} & \num{6.70(34)e+04} & \num{3.67(18)e+04} \\
	\hline
    \end{tabular}
    \caption{Halbwertszeiten \cite{521versuchsanleitung} der Isotope sowie Aktivitäten der Proben}
    \label{tab:activity}
\end{table}


Mit diesen Ausdrücken, und $p_\text{Übergang}$ der Wahrscheinlichkeit eines Übergangs, kann die Anzahl der Photonen die den Detektor erreichen geschrieben werden als
\begin{align}
	N_{E_\gamma, \text{Detektor}} = A(t) \cdot T_\text{Messung} \cdot q = A_0 e^{-\frac{t \ln(2)}{T_{1/2}}} T_\text{Messung} \frac12 \bigg (1 - \sqrt{1 + (\frac{r_\text{Detektor}}{d_\text{Probe}})}^{-1}\bigg) \cdot p_\text{Übergang}
\end{align}


Die Anzahl der im Phtotopeak registrierten Photonen ist direkt als Anpassungsparameter in Tabelle \todo{refrence} verfügbar.

\todo{Tabelle mit beliebig vielen zwischenwerten}

\section{Vergleich der Detektortypen}
Mit diesen Überlegungen und $p_\text{Übergang} = \SI{94.7(2)}{\percent}$ \cite{521versuchsanleitung} wird für die $\SI{661.7}{\kilo\electronvolt}$ Linie von Cäsium\-137 Quelle die Effizienz in beiden Detektoren bestimmt.
Es ergibt sich eine Effizienz des HPGE Detektors von $\epsilon_\text{HPGE} = \num{0.0567(68)}$ ($\approx 5.7\%$) und für NaI $\epsilon_\text{NaI} = 0.101(12)$ ($\approx 10.1\%$).

Man erkennt bei dieser Linie eine deutlich höhere Effizienz des Szintillationsdetektors, was genau den Erwartungen entspricht \todo{cite}.


\section{Effizienzkurve HPGE}

Nun wird für alle in \todo{refrence chapter} identifizierten Europium Linien die Effizienzen im HPGE Detektor bestimmt. Die Wahrscheinlichkeiten $p_\text{Übergang}$ sind hierbei der Praktikumsanleitung entnommen.
Die so gewonnenen Effizienzen werden in Abbildung \ref{fig:ge_efficiency} aufgetragen.
Die Energien werden auf $\SI{100}{\kilo\electronvolt}$ normiert und logarithmiert und gegen die logarithmierten Effizienzen Aufgetragen. An diese Werte wird nun ein Polynom zweiter Ordnung angepasst %\cite{}.

\jafps{ge_efficiency}{Effizienzkurve des HPGE Detektors mit Effizienz für Cäsium eingezeichnet}{0.75}
